\section{Coordinate sections and unrestricted coefficient growth}\label{sec:universality}

The bounded-rank hypothesis cannot be removed from the coefficient bounds.
To show this, we transfer transportation universality to the original
coordinates of a sparse bilinear hull. The universality theorem itself is
classical: \citet[Theorem~1.1]{DeLoeraOnn2006} represent every bounded
rational standard-form polytope by a slim three-way transportation
polytope, with a coordinate-erasing bijection. Their construction
(Section~3.3, printed page~816) places the retained coordinates in the
first layer through the injection $(i,j,k)\mapsto((i,j),(1,k),1)$.
Composing the coordinate injections of the construction therefore keeps
any specified original coordinates in that layer. Their Theorem~1.2 and
the discussion on pages~807--808 already establish bitransportation
universality and its two-commodity network interpretation.
\Cref{thm:universality} concerns only the further transfer to a
\emph{sparse bilinear coordinate section}, not universality of
multicommodity networks.

\begin{theorem}[Sparse coordinate-section transfer]\label{thm:universality}
Let $Q\subseteq\R_+^p$ be a nonempty bounded rational polytope, given by a
rational inequality description. In polynomial time one can construct an
acyclic four-layer network, a sparse observation set with two explicit
simplex variables, and a coordinate affine subspace fixing every original hull
coordinate except $p$ observed products, such that the hull section in
these $p$ coordinates is exactly $Q$. All arcs have the same positive
integer capacity $B$, and the flow value is $B$.
After dividing every flow and product coordinate by $B$, all capacities
and the flow value are one and the section becomes $Q/B$.
\end{theorem}
\begin{proof}
Introduce nonnegative slack variables for the inequalities of $Q$ and
clear equation denominators. Since $Q$ is bounded, its standard-form
extension is bounded as well. Apply the preceding transportation theorem
and its first-layer coordinate injection. Write the resulting table
polytope as
\begin{equation}\label{eq:universal-table}
 \mathcal T=\left\{t\ge0:
 \sum_{k=1}^3t_{ijk}=U_{ij},\quad
 \sum_jt_{ijk}=V_{ik},\quad
 \sum_it_{ijk}=W_{jk}\right\},
\end{equation}
where $i\in[r]$, $j\in[c]$, and all margins are nonnegative integers.
There are $p$ distinct designated cells $(i,j,1)$ whose values range
exactly over $Q$; erasing the other cells includes erasing the slack
coordinates. Nonemptiness implies consistency of the margins.

We may require $D_k:=\sum_iV_{ik}=\sum_jW_{jk}>0$ for every layer.
To arrange this, append one row and one column, put all old--new cross-cell
two-dimensional margins equal to zero, and put the new corner margin
equal to three. Assign new row and column layer margins equal to one
in each layer. Nonnegativity forces all cross entries to zero and each
corner entry to one. This padding does not change any old cell and
increases each $D_k$ by one. In the following, $r,c$ include the padding.

Put $B=D_1+D_2+D_3$ and build the network
\[
 s\longrightarrow r_i\longrightarrow c_j\longrightarrow t
 \qquad(i\in[r],\ j\in[c]),
\]
with every indicated arc present, every capacity $B$, and net supply
$B$ at $s$ and demand $B$ at $t$. Observe labels 1 and 2 on every
boundary arc $sr_i,c_jt$. On interior arcs observe label 1 only at
the $p$ designated cells. Thus there are $rc+r+c$ arcs and
$2(r+c)+p$ observed products. Fix the coordinates
\begin{align*}
 x_{r_ic_j}&=U_{ij},&
 x_{sr_i}&=\sum_kV_{ik},&x_{c_jt}&=\sum_kW_{jk},\\
 y_k&=D_k/B,&
 z_{sr_i,k}&=V_{ik},&z_{c_jt,k}&=W_{jk}\quad(k=1,2).
\end{align*}
These fix every original coordinate other than the designated interior
products. The residual weight is $D_3/B$.

By \cref{prop:disaggregation}, a point of this section has three state
flows. Their interior entries give a table satisfying
\eqref{eq:universal-table}: the aggregate interior flows give $U$,
the two observed boundary layers give $V,W$ for $k=1,2$, and
subtraction from the aggregate gives the third layer. Conversely any
table supplies the three state flows, including their boundary arcs.
Every arc of a nonnegative flow of value $D_k$ on this acyclic network
has flow at most $D_k$; for example, decompose it into directed
source--sink paths by repeatedly subtracting a positive path flow.
Hence the scaled capacity bounds $f^k_e\le(D_k/B)B$ hold.
The designated products are precisely the designated table entries.
This proves exactness directly; it does not rely on forming a hull that
retains all interior products and then arguing that its facets survive
projection.

All transformations have polynomial output size and rational encoding
length; the padding adds only one row and one column. Finally, the
invertible change $x'=x/B$, $z'=z/B$, $y'=y$ converts the network
to unit capacities and unit flow value, and scales the free section
coordinates by $1/B$.
\end{proof}

\begin{corollary}[Large coefficients with two simplex variables]\label{cor:universal-coefficients}
For every integer $M\ge2$ there is a sparse bilinear hull on an acyclic
unit-capacity unit-flow network, with $m=2$, whose every finite
original-coordinate inequality description contains an inequality with
two observed-product coefficients in ratio $M$. This ratio is unchanged
by adding affine-hull equations. The instances can be chosen to have
encoding length polynomial in $\log M$, with all numerical data in the
original nonlinear model in $\{-1,0,1\}$. In particular, the magnitudes
required in integer inequality descriptions are not polynomially bounded
by the original input length.
\end{corollary}
\begin{proof}
Apply \cref{thm:universality} to
$Q_M=\{(U,V)\ge0:U+MV\le1\}$. After normalization the free
coordinate section is $\{(U,V)\ge0:U+MV\le1/B\}$.
At an interior point of its sloping edge the section is locally exactly
one half-plane. \Cref{lem:section-facet} gives the ratio $M$, including
invariance under affine-hull equations. Two nonzero integer coefficients
in this ratio have maximum magnitude at least $M$.
The transportation construction has polynomial size in the encoding
length of $Q_M$, which is $O(\log M)$; normalization leaves only the
incidence, unit capacities, unit balances, simplex constraints, and
selected bilinear equations in the original model. The rational constants
fixing the section are not additional data of that model.
Taking $M=2^k$ proves the last assertion.
\end{proof}

These hulls are $0/1$ polytopes. A nonnegative unit flow on an acyclic
network is a convex combination of path incidence vectors. Pairing each
path with each simplex vertex gives $0/1$ values of $x,y,z$, and by
\cref{prop:disaggregation} their convex hull is exactly the sparse hull.
Fractional coordinate sections of a $0/1$ polytope need not be integral.
Large coefficients for general $0/1$ polytopes and ill-conditioned $0/1$
matrices are well established; see, for example, \citet{AlonVu1997}.
The restriction here is to selected flow--simplex products with only two
explicit states. The lower bound concerns coefficient \emph{magnitudes};
their bit lengths can remain polynomial. It does not imply separation
hardness, large extension complexity, or a corresponding statement for
planar networks.
